% ============================================================================
% Chapter 2, Section 2: Motivation & Context
% ============================================================================

\subsection{Motivation \& Context}\label{sec:motivation}

\subsubsection{The Long-Context Setting}\label{sec:long_context}

Benchmarks have historically served as catalysts for advancing AI capabilities. The Arcade Learning Environment (ALE)~\citep{bellemare2013arcade} catalyzed a decade of reinforcement learning progress by providing a standardized suite of Atari games~\citep{mnih2015human}. MineRL~\citep{guss2019minerl} extended this paradigm to open-ended environments with long-horizon objectives in Minecraft. More recently, NeurIPS competitions such as Neural MMO~\citep{liu2023neurips, suarez2023neural} and Lux AI~\citep{tao2024lux} have pushed the frontier on multi-agent cooperation and resource management. In each case, the availability of a standardized evaluation framework preceded---and arguably enabled---rapid methodological progress.

The current generation of frontier AI systems operates increasingly in long-horizon, autonomous settings. Anthropic's extended thinking capabilities enable Claude to reason over tens of thousands of actions~\citep{anthropic2025visible}. Google's Gemini 2.5 technical report documents agents that maintain coherent behavior across hundreds of hours of continuous interaction~\citep{gemini2p5report}. Commercial coding agents such as Claude Code~\citep{anthropic2025claudecode} and OpenHands~\citep{wang2024openhands} navigate multi-file codebases over extended sessions, performing planning, tool use, and context management at scales that would have been infeasible a year ago. These developments raise the importance of having benchmarks that natively test in the long-context, embodied setting---where agents must sustain coherent reasoning across thousands of sequential decisions, manage growing context windows, and integrate multimodal perception with strategic planning.

Yet existing benchmarks largely fail to stress these capabilities simultaneously. Standard LLM evaluations focus on short-context question answering, coding, or mathematical reasoning. Game benchmarks tend to isolate individual axes: imperfect-information games emphasize short-episode strategic reasoning~\citep{brown2018superhuman, brown2019superhuman}, while open-world environments like Minecraft test exploration without adversarial pressure. Recent orthogonality analysis via BenchPress~\citep{papailiopoulos2026benchpress} has confirmed that Pok\'{e}mon gameplay captures capabilities not predicted by the 49 benchmarks in their evaluation matrix---suggesting that this domain tests a genuinely distinct and underserved axis of model competence.

\subsubsection{LLM Agents in Pok\'{e}mon Games}\label{sec:pokemon_agents}

Pok\'{e}mon has emerged as a prominent testbed for evaluating frontier AI systems, driven by a convergence of properties that make it uniquely challenging. The franchise's RPG titles---particularly \textit{Pok\'{e}mon Emerald}, \textit{Red}, and \textit{Crystal}---feature expansive, partially observable worlds that require effective long-horizon planning, strategy adaptation, and knowledge integration to progress through a structured narrative. With approximately $10^{564}$ possible battle states~\citep{karten2025pokeagent}, team building across 1,000+ species with diverse movesets and abilities, and a competitive metagame that evolves continuously, the games present a complexity and dynamism that exceeds most existing benchmarks.

Several high-profile projects have demonstrated both the appeal and the difficulty of this domain. Claude Plays Pok\'{e}mon, a public Twitch stream, demonstrated extended thinking capabilities over approximately 35,000 actions to complete a small section of Pok\'{e}mon Red, but exposed struggles with basic spatial reasoning, vision interpretation, and obstacle navigation~\citep{engadget2025claudeplayspokemon, anthropic2025visible}. The Gemini Plays Pok\'{e}mon (GPP) project used Google's Gemini 2.5 Pro to play through Pok\'{e}mon Blue, ultimately completing the game in approximately 406 hours of cumulative gameplay, though the system frequently entered repetitive loops as the context window grew beyond 100K tokens~\citep{zhang2025geminiplayspokemon, gemini2p5report}. Later GPP runs with Gemini 3 Pro demonstrated substantial improvements: completing Pok\'{e}mon Crystal without losing a single battle, using half the turns and 60\% fewer tokens to reach the same milestones as its predecessor~\citep{zhang2025gemini3playspokemon}. OpenAI's GPT-5 completed Pok\'{e}mon Red in 6,470 steps~\citep{engadget2025gptplayspokemon}, establishing a speed benchmark that nonetheless remained difficult to compare against other systems due to differences in evaluation methodology.

These demonstrations powerfully validated Pok\'{e}mon as a testbed for embodied AI, but the efforts were fundamentally fragmented. Different projects used different games (Red, Blue, Crystal, Emerald), different harnesses with varying levels of built-in assistance, and different evaluation criteria. Was Gemini 3 Pro's completion of Crystal comparable to Claude's progress in Red? Did GPT-5's step count account for the same mechanics as the GPP runs? By conflating harness engineering with model capability, it became impossible to attribute success to the agent architecture, the underlying model, or hardcoded assumptions that simplified perception and navigation. The importance of standardized evaluation in games AI is well established: ALE catalyzed a decade of RL progress, while MineRL demonstrated the value of shared benchmarks for open-ended environments~\citep{milani2023solvingfuzzytaskshuman, shah2022retrospective}. Pok\'{e}mon demands---and now receives---similar standardization through the \pokeagent{} Challenge.

Pok\'{e}mon also offers a distinctive form of out-of-distribution evaluation. While extensive Pok\'{e}mon knowledge exists in pretraining corpora---move types, damage formulas, competitive tier lists---translating that latent knowledge into effective multi-turn sequential decision-making under partial observability is fundamentally different from the recognition and retrieval tasks where data contamination typically inflates performance. The competitive metagame shifts continuously as players develop new strategies and governing bodies rebalance tiers, creating natural distribution shifts that test an agent's ability to adapt rather than memorize.

\subsubsection{Game AI Benchmarks}\label{sec:game_benchmarks}

The need for standardization in Pok\'{e}mon-playing AI reflects a broader pattern in the history of game AI benchmarks. Game environments have driven major advances across the field: superhuman board game play through self-play reinforcement learning~\citep{silver2018general, silver2017mastering}, superhuman poker through equilibrium-finding algorithms~\citep{brown2018superhuman, brown2019superhuman}, grandmaster-level StarCraft~II through multi-agent RL~\citep{vinyals2019grandmaster}, and human-level Diplomacy through the integration of language models with strategic reasoning~\citep{meta2022cicero}. In each case, the benchmark both measured and drove progress.

A consistent pattern emerges across these results: RL achieves superhuman performance in fully observable, deterministic settings, but this margin erodes in stochastic, partially observable environments~\citep{samuel1959}. LLM-based agents consistently lag behind specialist RL and search systems in game settings~\citep{kaggle2025gamearena, paglieri2024balrog}, particularly when long-horizon planning and spatial reasoning are required. Hanabi~\citep{bard2020hanabi} and FightLadder~\citep{li2024fightladder} have advanced imperfect-information and competitive evaluation, while NetHack~\citep{kuttler2020nethack} and BALROG~\citep{paglieri2024balrog} benchmark long-horizon reasoning for both RL and LLM agents.

However, none of these benchmarks combines adversarial reasoning with long-horizon planning at the scale of a living competitive ecosystem, and most rely on symbolic state representations rather than visual perception. Pok\'{e}mon offers a unique combination: an enormous partially observed state space, a visual RPG requiring pixel-level perception, the need for sustained multi-hour reasoning, and a massive active player base generating continuous human data and an evolving competitive metagame. The NeurIPS 2025 \pokeagent{} Challenge confirmed this benchmark's difficulty and value: over 100 teams submitted solutions, 650+ researchers joined the competition Discord, and the results revealed a clear capability hierarchy---specialist RL and search methods dominated LLM approaches in competitive battling, while no raw frontier model achieved non-trivial progress in speedrunning without a sophisticated harness~\citep{karten2025pokeagent}. These gaps remain far from closed and motivate the benchmark, harness, and algorithmic contributions of this thesis.
